% --------------------------------------------------
% Chapter 04
% MLOps Architecture
% --------------------------------------------------

\label{chap:mlops-architecture}

\section{Learning Objectives}

After completing this chapter, learners will be able to:

\begin{itemize}
    \item Understand the major components of an MLOps architecture.
    \item Explain the role of each architectural layer.
    \item Describe data, model, and deployment flows.
    \item Understand how MLOps principles are implemented in practice.
    \item Design an enterprise-grade MLOps platform.
\end{itemize}

\section{Introduction}

MLOps principles provide the foundation for operational machine learning.

However, principles alone are insufficient.

Organizations must implement technical architectures capable of supporting:

\begin{itemize}
    \item Data management
    \item Model development
    \item Experiment tracking
    \item Deployment
    \item Monitoring
    \item Governance
\end{itemize}

This chapter introduces a reference MLOps architecture commonly used in modern machine learning platforms.

\section{Why Architecture Matters}

Many machine learning projects begin with isolated notebooks and local experimentation.

While this approach is useful for prototyping, it quickly becomes problematic when organizations need to:

\begin{itemize}
    \item Support multiple users
    \item Manage hundreds of models
    \item Deploy models at scale
    \item Satisfy regulatory requirements
\end{itemize}

Without architecture, organizations often face:

\begin{itemize}
    \item Duplicate work
    \item Data inconsistencies
    \item Deployment failures
    \item Governance deficiencies
\end{itemize}

A well-designed architecture provides standardization, scalability, and maintainability.

\section{Reference Architecture Overview}

A typical enterprise MLOps platform contains the following layers:

\begin{center}
Data Sources
\end{center}

\begin{center}
$\downarrow$
\end{center}

\begin{center}
Data Pipeline Layer
\end{center}

\begin{center}
$\downarrow$
\end{center}

\begin{center}
Feature Store
\end{center}

\begin{center}
$\downarrow$
\end{center}

\begin{center}
Training Pipeline
\end{center}

\begin{center}
$\downarrow$
\end{center}

\begin{center}
Experiment Tracking
\end{center}

\begin{center}
$\downarrow$
\end{center}

\begin{center}
Model Registry
\end{center}

\begin{center}
$\downarrow$
\end{center}

\begin{center}
Deployment Layer
\end{center}

\begin{center}
$\downarrow$
\end{center}

\begin{center}
Monitoring Layer
\end{center}

\begin{center}
$\downarrow$
\end{center}

\begin{center}
Governance Layer
\end{center}

\section{Data Sources Layer}

\subsection{Purpose}

The Data Sources Layer provides the raw information required to train and operate machine learning models.

This layer is often the most diverse component of the architecture.

\subsection{Examples of Data Sources}

Typical sources include:

\begin{itemize}
    \item Relational databases
    \item Data warehouses
    \item Data lakes
    \item CRM systems
    \item ERP systems
    \item Transaction platforms
    \item External vendors
    \item Public datasets
\end{itemize}

\subsection{Credit Scoring Example}

A credit risk model may use information from:

\begin{itemize}
    \item Loan applications
    \item Customer profiles
    \item Payment histories
    \item Credit bureau reports
    \item Macroeconomic indicators
\end{itemize}

\subsection{Architectural Considerations}

Organizations must address:

\begin{itemize}
    \item Data quality
    \item Data ownership
    \item Data lineage
    \item Data governance
    \item Security controls
\end{itemize}

\section{Data Pipeline Layer}

\subsection{Purpose}

The Data Pipeline Layer transforms raw data into usable analytical datasets.

This layer performs all required preparation activities before model training.

\subsection{Typical Activities}

Examples include:

\begin{itemize}
    \item Data ingestion
    \item Data cleansing
    \item Schema validation
    \item Data integration
    \item Data transformation
    \item Quality controls
\end{itemize}

\subsection{Automation}

Data pipelines should be automated whenever possible.

Automation improves:

\begin{itemize}
    \item Reliability
    \item Repeatability
    \item Scalability
\end{itemize}

\subsection{Common Technologies}

Examples include:

\begin{itemize}
    \item Apache Airflow
    \item Azure Data Factory
    \item Databricks
    \item Spark
    \item Prefect
\end{itemize}

\section{Feature Store}

\subsection{Purpose}

Feature engineering is one of the most critical activities in machine learning.

A Feature Store provides a centralized repository for storing, managing, and serving features used by machine learning models.

The objective is to ensure consistency between training and inference environments.

\subsection{Why Feature Stores Exist}

Without a Feature Store, organizations frequently encounter:

\begin{itemize}
    \item Duplicate feature calculations
    \item Inconsistent definitions
    \item Training-serving skew
    \item Increased maintenance costs
\end{itemize}

Different teams may compute the same feature in different ways, leading to inconsistent model behavior.

\subsection{Typical Responsibilities}

A Feature Store typically provides:

\begin{itemize}
    \item Feature storage
    \item Feature versioning
    \item Metadata management
    \item Online serving
    \item Offline serving
\end{itemize}

\subsection{Offline Features}

Offline features are used during model development and training.

Characteristics include:

\begin{itemize}
    \item Large datasets
    \item Historical information
    \item Batch processing
\end{itemize}

\subsection{Online Features}

Online features support real-time inference.

Characteristics include:

\begin{itemize}
    \item Low latency
    \item Near real-time updates
    \item API accessibility
\end{itemize}

\subsection{Feature Metadata}

Each feature should be documented.

Typical metadata includes:

\begin{itemize}
    \item Feature name
    \item Description
    \item Data type
    \item Owner
    \item Refresh frequency
    \item Source systems
\end{itemize}

\subsection{Credit Scoring Example}

Examples of features include:

\begin{itemize}
    \item Debt-to-income ratio
    \item Credit utilization rate
    \item Number of recent inquiries
    \item Average account age
    \item Number of missed payments
\end{itemize}

These features should be calculated once and reused consistently across all models.

\subsection{Popular Technologies}

Examples include:

\begin{itemize}
    \item Feast
    \item Databricks Feature Store
    \item Tecton
    \item Azure Feature Store
\end{itemize}

\section{Training Pipeline Layer}

\subsection{Purpose}

The Training Pipeline Layer automates model development activities.

Instead of manually executing notebooks and scripts, organizations create standardized training workflows.

\subsection{Typical Pipeline Stages}

A training pipeline may include:

\begin{enumerate}
    \item Dataset retrieval
    \item Data validation
    \item Feature generation
    \item Model training
    \item Hyperparameter tuning
    \item Model evaluation
    \item Artifact generation
\end{enumerate}

\subsection{Benefits}

Training pipelines improve:

\begin{itemize}
    \item Reproducibility
    \item Scalability
    \item Reliability
    \item Productivity
\end{itemize}

\subsection{Pipeline Inputs}

Typical inputs include:

\begin{itemize}
    \item Dataset version
    \item Feature definitions
    \item Model configuration
    \item Hyperparameters
\end{itemize}

\subsection{Pipeline Outputs}

Typical outputs include:

\begin{itemize}
    \item Trained model
    \item Metrics
    \item Logs
    \item Feature importance reports
    \item Validation reports
\end{itemize}

\subsection{Hyperparameter Optimization}

Many organizations integrate automated optimization techniques.

Examples include:

\begin{itemize}
    \item Grid Search
    \item Random Search
    \item Bayesian Optimization
\end{itemize}

\subsection{Distributed Training}

Large-scale machine learning systems may require distributed computing resources.

Benefits include:

\begin{itemize}
    \item Faster training
    \item Larger datasets
    \item More complex models
\end{itemize}

\subsection{Training Environment Standardization}

Training environments should be standardized through:

\begin{itemize}
    \item Containers
    \item Dependency management
    \item Infrastructure-as-Code
\end{itemize}

This ensures that models can be reproduced consistently.

\subsection{Training Pipeline Technologies}

Examples include:

\begin{itemize}
    \item Azure ML Pipelines
    \item Kubeflow Pipelines
    \item Airflow
    \item Vertex AI Pipelines
    \item Databricks Workflows
\end{itemize}

\section{Experiment Tracking Layer}

\subsection{Purpose}

Machine learning development is an iterative and experimental process.

Data scientists frequently train multiple models before selecting a candidate for deployment.

The Experiment Tracking Layer records all relevant information associated with these experiments.

\subsection{Objectives}

Experiment tracking aims to:

\begin{itemize}
    \item Improve reproducibility
    \item Facilitate collaboration
    \item Support model comparison
    \item Enable auditability
\end{itemize}

\subsection{Tracked Information}

Typical experiment metadata includes:

\begin{itemize}
    \item Experiment identifier
    \item Dataset version
    \item Feature set version
    \item Training code version
    \item Hyperparameters
    \item Metrics
    \item Execution timestamp
    \item User information
\end{itemize}

\subsection{Metrics Tracking}

Examples of tracked metrics include:

\begin{itemize}
    \item Accuracy
    \item Precision
    \item Recall
    \item ROC-AUC
    \item F1 Score
    \item KS Statistic
\end{itemize}

\subsection{Artifact Tracking}

In addition to metrics, organizations should retain:

\begin{itemize}
    \item Trained models
    \item Validation reports
    \item Feature importance reports
    \item Training logs
    \item Configuration files
\end{itemize}

\subsection{Benefits}

Experiment tracking allows teams to answer questions such as:

\begin{itemize}
    \item Which model was deployed?
    \item Which dataset was used?
    \item Which hyperparameters produced the best results?
    \item Can the model be reproduced?
\end{itemize}

\subsection{Credit Scoring Example}

Suppose a team trains twenty candidate credit scoring models.

Experiment tracking enables comparison across:

\begin{itemize}
    \item Logistic Regression
    \item Random Forest
    \item XGBoost
    \item LightGBM
\end{itemize}

without losing historical information.

\subsection{Popular Technologies}

Examples include:

\begin{itemize}
    \item MLflow Tracking
    \item Azure Machine Learning
    \item Weights \& Biases
    \item Neptune
\end{itemize}

\section{Model Registry Layer}

\subsection{Purpose}

The Model Registry acts as the central repository for machine learning models.

It provides lifecycle management for model artifacts.

Every approved model should be stored in a registry before deployment.

\subsection{Why a Registry Is Necessary}

Without a registry, organizations often struggle to identify:

\begin{itemize}
    \item The current production model
    \item Previous versions
    \item Validation status
    \item Approval history
\end{itemize}

This creates operational and governance risks.

\subsection{Model Lifecycle}

A typical lifecycle includes:

\begin{center}
Development
$\rightarrow$
Validation
$\rightarrow$
Approval
$\rightarrow$
Production
$\rightarrow$
Retirement
\end{center}

\subsection{Model Metadata}

A registry should maintain:

\begin{itemize}
    \item Model name
    \item Version number
    \item Training dataset reference
    \item Feature set reference
    \item Metrics
    \item Approval status
    \item Deployment history
\end{itemize}

\subsection{Version Management}

Example:

\begin{itemize}
    \item Credit Model v1.0
    \item Credit Model v1.1
    \item Credit Model v2.0
\end{itemize}

Each version should remain accessible for future audits.

\subsection{Approval Workflows}

Organizations typically define approval processes involving:

\begin{itemize}
    \item Data Scientists
    \item Model Validators
    \item Business Owners
    \item Risk Managers
\end{itemize}

Only approved models should progress to production.

\subsection{Model Promotion}

A registry supports controlled promotion between environments:

\begin{center}
Development
$\rightarrow$
Testing
$\rightarrow$
Staging
$\rightarrow$
Production
\end{center}

\subsection{Governance Benefits}

A registry improves:

\begin{itemize}
    \item Traceability
    \item Auditability
    \item Compliance
    \item Operational control
\end{itemize}

\subsection{MLflow Registry}

MLflow Model Registry provides:

\begin{itemize}
    \item Model versioning
    \item Stage transitions
    \item Metadata management
    \item Approval workflows
\end{itemize}

\subsection{Azure ML Registry}

Azure Machine Learning extends registry capabilities through:

\begin{itemize}
    \item Centralized model management
    \item Enterprise governance
    \item Security controls
    \item Multi-environment deployment support
\end{itemize}

\subsection{Best Practices}

Organizations should:

\begin{itemize}
    \item Register all production models
    \item Maintain version history
    \item Record approval decisions
    \item Link models to validation reports
    \item Document deployment decisions
\end{itemize}

\section{Deployment Layer}

\subsection{Purpose}

The Deployment Layer is responsible for making trained machine learning models available to users, applications, and business processes.

Deployment transforms a validated model into an operational service capable of generating predictions.

The objective is to ensure that models can be:

\begin{itemize}
    \item Accessed reliably
    \item Scaled efficiently
    \item Updated safely
    \item Monitored continuously
\end{itemize}

\subsection{Deployment Challenges}

Deploying machine learning systems introduces several challenges:

\begin{itemize}
    \item Dependency management
    \item Environment consistency
    \item Scalability requirements
    \item Security controls
    \item Model version management
\end{itemize}

A model that performs well during experimentation may fail in production if these challenges are not addressed.

\section{Containerization}

\subsection{Why Containers?}

Containers provide a standardized execution environment.

A container packages:

\begin{itemize}
    \item Application code
    \item Dependencies
    \item Libraries
    \item Runtime configuration
\end{itemize}

This ensures consistency across development, testing, and production environments.

\subsection{Docker}

Docker has become the industry standard for machine learning deployment.

Benefits include:

\begin{itemize}
    \item Portability
    \item Reproducibility
    \item Scalability
    \item Isolation
\end{itemize}

\subsection{Typical Container Architecture}

\begin{center}
Model
$\rightarrow$
Prediction API
$\rightarrow$
Docker Container
$\rightarrow$
Deployment Platform
\end{center}

\section{Batch Inference}

\subsection{Definition}

Batch inference generates predictions for large collections of records at scheduled intervals.

Typical examples include:

\begin{itemize}
    \item Monthly credit scoring
    \item Customer segmentation
    \item Portfolio risk assessment
\end{itemize}

\subsection{Characteristics}

Batch inference typically involves:

\begin{itemize}
    \item High volumes
    \item Scheduled execution
    \item Lower latency requirements
\end{itemize}

\subsection{Advantages}

Advantages include:

\begin{itemize}
    \item Simpler architecture
    \item Lower operational complexity
    \item Easier resource planning
\end{itemize}

\section{Online Inference}

\subsection{Definition}

Online inference generates predictions in real time.

Predictions are produced immediately after receiving a request.

Examples include:

\begin{itemize}
    \item Credit application scoring
    \item Fraud detection
    \item Recommendation systems
\end{itemize}

\subsection{Characteristics}

Online systems require:

\begin{itemize}
    \item Low latency
    \item High availability
    \item Scalability
    \item Fault tolerance
\end{itemize}

\subsection{Typical Architecture}

\begin{center}
Client Application
$\rightarrow$
REST API
$\rightarrow$
Model Service
$\rightarrow$
Prediction
\end{center}

\section{Model Serving}

\subsection{Purpose}

Model serving refers to the process of exposing machine learning models through standardized interfaces.

Serving systems provide prediction services to external consumers.

\subsection{Serving Interfaces}

Common interfaces include:

\begin{itemize}
    \item REST APIs
    \item gRPC Services
    \item Batch Jobs
    \item Streaming Systems
\end{itemize}

\subsection{Prediction APIs}

Prediction APIs typically perform:

\begin{enumerate}
    \item Request validation
    \item Feature preparation
    \item Model execution
    \item Response generation
\end{enumerate}

\subsection{FastAPI}

FastAPI has become a popular framework for machine learning APIs.

Benefits include:

\begin{itemize}
    \item High performance
    \item Automatic documentation
    \item Python integration
    \item Simplicity
\end{itemize}

\section{Deployment Strategies}

\subsection{Blue-Green Deployment}

Two identical environments are maintained:

\begin{itemize}
    \item Current production environment
    \item New candidate environment
\end{itemize}

Traffic is switched only after validation.

\subsection{Canary Deployment}

The new model receives a small percentage of traffic.

Examples:

\begin{itemize}
    \item 5\% of requests
    \item 10\% of requests
    \item 25\% of requests
\end{itemize}

Traffic increases gradually if performance remains acceptable.

\subsection{Shadow Deployment}

The new model receives production traffic without affecting business decisions.

Predictions are compared with the current production model.

This approach reduces deployment risk.

\section{Cloud Deployment Platforms}

\subsection{Azure Machine Learning Endpoints}

Azure ML provides managed endpoints for model serving.

Capabilities include:

\begin{itemize}
    \item Managed infrastructure
    \item Autoscaling
    \item Monitoring
    \item Security controls
\end{itemize}

\subsection{Kubernetes}

Large organizations frequently deploy machine learning services using Kubernetes.

Benefits include:

\begin{itemize}
    \item Scalability
    \item High availability
    \item Resource management
    \item Container orchestration
\end{itemize}

\subsection{Serverless Inference}

Serverless platforms automatically provision resources when predictions are requested.

Advantages include:

\begin{itemize}
    \item Reduced infrastructure management
    \item Cost optimization
    \item Automatic scaling
\end{itemize}

\section{Monitoring Layer}

\subsection{Purpose}

The Monitoring Layer provides continuous visibility into the health, performance, and behavior of machine learning systems operating in production.

Monitoring is essential because machine learning models are exposed to changing environments and evolving data distributions.

Without monitoring, organizations may be unaware of deteriorating model performance until significant business losses occur.

\subsection{Monitoring Objectives}

The primary objectives are:

\begin{itemize}
    \item Detect model degradation
    \item Detect data drift
    \item Detect concept drift
    \item Monitor operational performance
    \item Support governance requirements
\end{itemize}

\subsection{Operational Monitoring}

Operational monitoring focuses on system performance.

Typical metrics include:

\begin{itemize}
    \item Request volume
    \item Latency
    \item Response time
    \item Error rates
    \item Availability
    \item Resource utilization
\end{itemize}

\subsection{Data Monitoring}

Data monitoring focuses on incoming prediction data.

Typical controls include:

\begin{itemize}
    \item Missing value rates
    \item Outlier detection
    \item Schema validation
    \item Distribution changes
    \item Data freshness
\end{itemize}

\subsection{Model Monitoring}

Model monitoring focuses on predictive quality.

Examples include:

\begin{itemize}
    \item Accuracy
    \item Precision
    \item Recall
    \item ROC-AUC
    \item Calibration
\end{itemize}

\subsection{Business Monitoring}

Business monitoring evaluates whether the model continues to generate value.

Examples include:

\begin{itemize}
    \item Default rates
    \item Fraud losses
    \item Customer retention
    \item Revenue impact
\end{itemize}

\subsection{Alerting Mechanisms}

Monitoring systems should trigger alerts when predefined thresholds are exceeded.

Examples include:

\begin{itemize}
    \item Significant drift detected
    \item Prediction volume anomalies
    \item Increased latency
    \item Performance degradation
\end{itemize}

\subsection{Monitoring Technologies}

Examples include:

\begin{itemize}
    \item Azure Monitor
    \item Prometheus
    \item Grafana
    \item Datadog
    \item Evidently AI
\end{itemize}

\section{Governance Layer}

\subsection{Purpose}

The Governance Layer ensures that machine learning systems remain compliant, auditable, and aligned with organizational policies.

Governance spans the entire model lifecycle.

\subsection{Governance Objectives}

Organizations should be able to answer:

\begin{itemize}
    \item Which model is currently deployed?
    \item Who approved the deployment?
    \item Which data was used?
    \item How was the model validated?
    \item What monitoring controls exist?
\end{itemize}

\subsection{Key Governance Components}

Typical governance components include:

\begin{itemize}
    \item Model documentation
    \item Validation reports
    \item Approval workflows
    \item Audit trails
    \item Change management processes
\end{itemize}

\subsection{Model Documentation}

Documentation should include:

\begin{itemize}
    \item Business purpose
    \item Data sources
    \item Methodology
    \item Assumptions
    \item Limitations
    \item Validation results
\end{itemize}

\subsection{Auditability}

Every significant action should be traceable.

Examples include:

\begin{itemize}
    \item Dataset creation
    \item Model training
    \item Validation activities
    \item Deployment decisions
    \item Production changes
\end{itemize}

\subsection{Regulatory Context}

Examples of relevant frameworks include:

\begin{itemize}
    \item SR 11-7
    \item ECB TRIM
    \item EBA Guidelines
    \item EU AI Act
    \item Internal Risk Policies
\end{itemize}

\section{Responsible AI Layer}

\subsection{Why Responsible AI Matters}

Modern organizations must ensure that machine learning systems are not only accurate but also trustworthy.

Responsible AI extends traditional MLOps by introducing additional controls focused on ethics and fairness.

\subsection{Core Principles}

Common Responsible AI principles include:

\begin{itemize}
    \item Fairness
    \item Transparency
    \item Accountability
    \item Reliability
    \item Privacy
    \item Security
\end{itemize}

\subsection{Fairness Assessment}

Organizations should evaluate whether model outcomes differ significantly across demographic groups.

Examples include:

\begin{itemize}
    \item Gender
    \item Age groups
    \item Geographic regions
\end{itemize}

\subsection{Explainability}

Stakeholders often require explanations for model decisions.

Common explainability techniques include:

\begin{itemize}
    \item Feature Importance
    \item SHAP Values
    \item LIME
    \item Partial Dependence Analysis
\end{itemize}

\subsection{Responsible AI Tooling}

Examples include:

\begin{itemize}
    \item Microsoft Responsible AI Dashboard
    \item Fairlearn
    \item SHAP
    \item InterpretML
\end{itemize}

\section{End-to-End Reference Architecture}

A mature MLOps platform can be represented as:

\begin{center}
Data Sources
\end{center}

\begin{center}
$\downarrow$
\end{center}

\begin{center}
Data Pipelines
\end{center}

\begin{center}
$\downarrow$
\end{center}

\begin{center}
Feature Store
\end{center}

\begin{center}
$\downarrow$
\end{center}

\begin{center}
Training Pipelines
\end{center}

\begin{center}
$\downarrow$
\end{center}

\begin{center}
Experiment Tracking
\end{center}

\begin{center}
$\downarrow$
\end{center}

\begin{center}
Model Registry
\end{center}

\begin{center}
$\downarrow$
\end{center}

\begin{center}
Deployment Layer
\end{center}

\begin{center}
$\downarrow$
\end{center}

\begin{center}
Monitoring Layer
\end{center}

\begin{center}
$\downarrow$
\end{center}

\begin{center}
Governance Layer
\end{center}

\begin{center}
$\downarrow$
\end{center}

\begin{center}
Responsible AI Controls
\end{center}

\section{Summary}

Enterprise MLOps architectures transform machine learning from isolated experimentation into scalable production systems.

A complete architecture typically includes:

\begin{itemize}
    \item Data Sources
    \item Data Pipelines
    \item Feature Store
    \item Training Pipelines
    \item Experiment Tracking
    \item Model Registry
    \item Deployment Services
    \item Monitoring Frameworks
    \item Governance Controls
    \item Responsible AI Mechanisms
\end{itemize}

Together, these components provide the foundation required to build reliable, scalable, auditable, and trustworthy machine learning systems.

The next chapter compares MLOps, DevOps, and DataOps and explains how these disciplines complement one another within modern organizations.